
The approach nevertheless has two disadvantages for deployment. First, each channel receives one threshold calibrated on the whole session, so the pipeline cannot run online as it stands and does not follow the changes in blink amplitude and duration that occur as sleepiness develops within a drive~\citep{akerstedt2010reaction,ingre2006subjective}. Re-estimating the threshold over a moving window of recent epochs could address both, and real-time threshold adaptation has been achieved for unsupervised blink detection~\citep{chang2016unsupervised}. Second, candidate epochs are selected by amplitude alone, so movement and electrode-contact artefacts common outside the laboratory~\citep{casson2019wearable} could enter the calibration pool. The median-based estimate limits the influence of isolated extremes, but frequent artefacts could still raise the threshold. Separately from these disadvantages, the comparison used the full montage and each method's best channel per session, so whether the margin holds for the few-channel frontal systems used in drowsiness monitoring~\citep{ko2020eyeblink,shahbakhti2022simultaneous} remains to be tested.
